\documentclass[12pt,a4paper,openany,oneside]{book}

% --- Paquetes ---
\usepackage[utf8]{inputenc}
\usepackage[spanish, es-noshorthands]{babel}
\usepackage{amsmath, amsfonts, amssymb}
\usepackage{graphicx}
\usepackage{geometry}
\usepackage{hyperref}
\usepackage{booktabs}
\usepackage{fancyhdr}
\usepackage{listings}
\usepackage{xcolor}
\usepackage{tikz}
\usepackage{pgfplots}
\usepackage{colortbl}
\usepackage{multirow}
\usepackage{hhline}
\usetikzlibrary{arrows.meta, positioning, shapes.geometric, calc, shapes}
\pgfplotsset{compat=1.18}

\geometry{margin=2.5cm}

% --- Colores y estilos para código Python ---
\definecolor{codegreen}{rgb}{0,0.6,0}
\definecolor{codegray}{rgb}{0.5,0.5,0.5}
\definecolor{codepurple}{rgb}{0.58,0,0.82}
\definecolor{backcolour}{rgb}{0.95,0.95,0.92}

\lstset{
    language=Python,
    backgroundcolor=\color{backcolour},   
    commentstyle=\color{codegreen},
    keywordstyle=\color{blue},
    numberstyle=\tiny\color{codegray},
    stringstyle=\color{codepurple},
    basicstyle=\ttfamily\small,
    breaklines=true,
    frame=single,
    showstringspaces=false
}

% --- Estilo de cabecera y pie ---
\pagestyle{fancy}
\fancyhf{}
\renewcommand{\headrulewidth}{0.4pt}
\renewcommand{\footrulewidth}{0.4pt}
\fancyhead[L]{\small Carlos César Sánchez Coronel}
\fancyhead[R]{\small Sesión 11: Modelos Complementarios}
\fancyfoot[L]{\small Carlos César Sánchez Coronel}
\fancyfoot[C]{\thepage}

% --- Portada ---
\title{
    \textbf{Sesión 11} \\ 
    \Huge \textbf{MODELOS COMPLEMENTARIOS} \\
    \Large Técnicas especializadas para problemas específicos
}
\author{
    \small Por \\ \vspace{0.2cm}
    \Large \textsc{Carlos César Sánchez Coronel}
}
\date{2026}

\begin{document}

\maketitle
\tableofcontents
\addcontentsline{toc}{chapter}{Índice general}

% ------------------------------------------------------------------
% CAPÍTULO 11: SESIÓN 11 - MODELOS COMPLEMENTARIOS
% ------------------------------------------------------------------
\chapter{Modelos Complementarios}

En las sesiones anteriores se han cubierto los modelos y técnicas más utilizados en machine learning. Sin embargo, en problemas reales surgen situaciones que requieren enfoques especializados: datos con outliers severos, relaciones no lineales que los árboles no capturan bien, necesidad de detectar clusters de forma arbitraria, o la obligación de interpretar jerarquías. Esta sesión presenta una colección de modelos complementarios que, aunque no son los más populares, resultan esenciales en contextos específicos y aparecen con frecuencia en entrevistas técnicas.

\section{Logro de la sesión}
Conocer otros algoritmos y técnicas que pueden aparecer en entrevistas o en problemas específicos, y entender cuándo aplicarlos según las características de los datos y los requisitos del problema.

\section{Regresión Polinómica}

\subsection{Extensión de la regresión lineal}
La regresión polinómica consiste en añadir términos polinómicos de las variables originales como nuevas características. Para una variable \(x\), el modelo de grado \(d\) es:
\[
y = \beta_0 + \beta_1 x + \beta_2 x^2 + \dots + \beta_d x^d + \varepsilon
\]
A pesar del nombre, sigue siendo un modelo lineal en los parámetros \(\beta\), por lo que puede estimarse con mínimos cuadrados ordinarios.

\subsection{Fundamento matemático}
La matriz de diseño se construye como:
\[
\mathbf{X} = \begin{pmatrix}
1 & x_1 & x_1^2 & \dots & x_1^d \\
1 & x_2 & x_2^2 & \dots & x_2^d \\
\vdots & \vdots & \vdots & \ddots & \vdots \\
1 & x_n & x_n^2 & \dots & x_n^d
\end{pmatrix}
\]
y se resuelve \(\hat{\beta} = (\mathbf{X}^T \mathbf{X})^{-1} \mathbf{X}^T \mathbf{y}\).

\subsection{Riesgo de overfitting}
A medida que aumenta el grado \(d\), el modelo se vuelve más flexible y puede ajustarse excesivamente a los datos de entrenamiento, especialmente en las regiones extremas. Es fundamental validar el grado adecuado mediante validación cruzada.

\begin{center}
\begin{tikzpicture}
\begin{axis}[
    xlabel={\(x\)},
    ylabel={\(y\)},
    xmin=-3, xmax=3,
    ymin=-5, ymax=15,
    grid=both,
    width=0.6\textwidth,
    legend pos=north west
]
\addplot[only marks, mark=*] coordinates {
    (-2.5, 2) (-2, 1) (-1.5, 0.5) (-1, 1) (-0.5, 2) (0, 3) (0.5, 5) (1, 7) (1.5, 10) (2, 12)
};
\addplot[blue, thick, domain=-2.5:2, samples=100] {0.5 + 1.2*x + 0.8*x^2};
\addplot[red, thick, domain=-2.5:2, samples=100] {0.5 + 1.2*x + 0.8*x^2 - 0.1*x^3 + 0.5*x^4 - 0.2*x^5};
\legend{Datos, Grado 2, Grado 5}
\end{axis}
\end{tikzpicture}
\caption{Regresión polinómica: grado 2 (azul) generaliza bien; grado 5 (rojo) sobreajusta.}
\end{center}

\subsection{Cuándo usar regresión polinómica}
\begin{itemize}
    \item Cuando se sospecha una relación no lineal pero suave.
    \item Como alternativa a modelos más complejos cuando se necesita interpretabilidad (los coeficientes de los términos polinómicos indican la forma de la relación).
    \item En problemas con pocas variables independientes.
\end{itemize}

\section{Regresión Robusta}

Los modelos de regresión estándar (mínimos cuadrados) son muy sensibles a outliers. La regresión robusta utiliza funciones de pérdida que penalizan menos los errores grandes.

\subsection{Regresión de Huber}
La función de pérdida de Huber combina el error cuadrático para errores pequeños con el error absoluto para errores grandes:
\[
L_{\delta}(r) = 
\begin{cases}
\frac{1}{2} r^2 & \text{si } |r| \le \delta \\
\delta(|r| - \frac{1}{2}\delta) & \text{en otro caso}
\end{cases}
\]
donde \(r = y - \hat{y}\) es el residuo y \(\delta\) es un umbral que controla la transición. Para residuos pequeños, se comporta como MSE (diferenciable y eficiente); para residuos grandes, como MAE (menos sensible a outliers).

\subsection{RANSAC (RANdom SAmple Consensus)}
RANSAC es un algoritmo iterativo para estimar parámetros de un modelo en presencia de un gran número de outliers. Su filosofía es diferente: no busca minimizar una función de pérdida sobre todos los datos, sino encontrar un conjunto de puntos "inliers" que sean consistentes con un modelo.

\begin{enumerate}
    \item Seleccionar aleatoriamente un subconjunto mínimo de puntos para ajustar un modelo (ej. 2 puntos para una recta).
    \item Calcular cuántos puntos del conjunto total están dentro de un umbral de distancia al modelo (inliers).
    \item Si el número de inliers supera un umbral, reajustar el modelo con todos los inliers.
    \item Repetir \(N\) iteraciones y quedarse con el modelo con más inliers.
\end{enumerate}

\begin{itemize}
    \item \textbf{Ventaja}: extremadamente robusto, puede tolerar más del 50\% de outliers.
    \item \textbf{Desventaja}: requiere ajustar parámetros (umbral de distancia, número de iteraciones) y no es determinista.
\end{itemize}

\subsection{Caso de uso: estimación de tendencia con outliers}
En datos de precios de viviendas, pueden existir propiedades atípicas (mansiones) que distorsionan la regresión lineal. Huber o RANSAC proporcionan una estimación más representativa de la relación para la mayoría de los datos.

\subsection{Comparación de pérdidas}
\begin{table}[h]
\centering
\begin{tabular}{|l|c|c|}
\hline
\textbf{Función} & \textbf{Ecuación} & \textbf{Sensibilidad a outliers} \\
\hline
MSE & \(r^2\) & Alta \\
MAE & \(|r|\) & Moderada \\
Huber & \(L_{\delta}(r)\) & Baja (con \(\delta\) adecuado) \\
\hline
\end{tabular}
\caption{Comparación de funciones de pérdida.}
\end{table}

\section{Support Vector Regression (SVR)}

SVR extiende las ideas de SVM a problemas de regresión. En lugar de minimizar el error cuadrático, busca una función \(f(x)\) que tenga al menos \(\varepsilon\) de desviación respecto a los valores reales, y al mismo tiempo sea lo más plana posible.

\subsection{Formulación}
Para un modelo lineal \(f(x) = \mathbf{w}^T x + b\), SVR resuelve:
\[
\min_{\mathbf{w}, b} \frac{1}{2} \|\mathbf{w}\|^2 + C \sum_{i=1}^n (\xi_i + \xi_i^*)
\]
sujeto a:
\[
y_i - (\mathbf{w}^T x_i + b) \le \varepsilon + \xi_i
\]
\[
(\mathbf{w}^T x_i + b) - y_i \le \varepsilon + \xi_i^*
\]
\[
\xi_i, \xi_i^* \ge 0
\]
donde:
\begin{itemize}
    \item \(\varepsilon\): margen de tolerancia (errores menores a \(\varepsilon\) no se penalizan).
    \item \(C\): parámetro de regularización que controla la penalización por errores mayores a \(\varepsilon\).
    \item \(\xi_i, \xi_i^*\): variables de holgura que miden la desviación fuera de la banda.
\end{itemize}

\subsection{Truco del kernel}
Al igual que en SVM clasificación, SVR puede usar kernels para modelar relaciones no lineales:
\[
f(x) = \sum_{i=1}^n (\alpha_i - \alpha_i^*) K(x_i, x) + b
\]
donde \(K\) es una función kernel (lineal, polinomial, RBF).

\subsection{Interpretación de parámetros}
\begin{itemize}
    \item \(\varepsilon\): define una banda "insensible" alrededor de la función. Un \(\varepsilon\) grande produce una función más plana, ignorando pequeñas desviaciones.
    \item \(C\): controla la penalización por puntos fuera de la banda. Valores grandes de \(C\) ajustan más los datos, pudiendo sobreajustar.
\end{itemize}

\subsection{Caso de uso: predicción de series temporales con outliers}
SVR es robusto a outliers si se elige un \(\varepsilon\) adecuado. Además, con kernel RBF puede capturar patrones no lineales.

\section{DBSCAN a fondo}

En la sesión 9 se introdujo DBSCAN. Aquí se profundiza en sus detalles y variantes.

\subsection{Elección de parámetros}
\begin{itemize}
    \item \textbf{eps (\(\varepsilon\))}: se puede estimar mediante el gráfico de distancia al k-ésimo vecino. Para un \(k\) dado (usualmente minPts), se ordenan las distancias y se busca un punto de inflexión (codo).
    \item \textbf{minPts}: regla general, minPts \(\ge\) dimensionalidad + 1. Valores típicos: 3-10.
\end{itemize}

\subsection{Métricas de distancia}
DBSCAN funciona con cualquier métrica. En espacios de alta dimensionalidad, la distancia euclidiana pierde significado; se pueden usar distancias basadas en coseno o Manhattan.

\subsection{Variantes}
\begin{itemize}
    \item \textbf{OPTICS} (Ordering Points To Identify the Clustering Structure): extiende DBSCAN para detectar clusters de densidad variable. Produce un diagrama de alcanzabilidad que permite visualizar la estructura jerárquica.
    \item \textbf{HDBSCAN}: variante jerárquica que encuentra clusters de diferentes densidades y es menos sensible a los parámetros.
\end{itemize}

\subsection{Ventajas sobre K-Means}
\begin{itemize}
    \item No asume clusters esféricos.
    \item No requiere especificar \(k\).
    \item Identifica puntos como ruido.
\end{itemize}

\subsection{Limitaciones}
\begin{itemize}
    \item Sensible a la elección de eps y minPts.
    \item Dificultad con densidades muy variables (aunque OPTICS lo mejora).
\end{itemize}

\section{Clustering Jerárquico a fondo}

El clustering jerárquico construye una jerarquía de clusters. Se profundiza en los tipos de enlace y la interpretación del dendrograma.

\subsection{Tipos de enlace y su efecto}
\begin{itemize}
    \item \textbf{Single linkage}: distancia mínima entre puntos de los clusters. Tiende a formar cadenas (efecto "encadenamiento").
    \item \textbf{Complete linkage}: distancia máxima. Produce clusters compactos y esféricos.
    \item \textbf{Average linkage}: promedio de distancias. Compromiso entre los anteriores.
    \item \textbf{Ward linkage}: minimiza el incremento de la suma de cuadrados intra-cluster. Tiende a producir clusters de tamaño similar.
\end{itemize}

\begin{center}
\begin{tikzpicture}
\begin{axis}[
    xlabel={\(x_1\)},
    ylabel={\(x_2\)},
    xmin=0, xmax=10,
    ymin=0, ymax=10,
    width=0.8\textwidth,
    title={Efecto del tipo de enlace}
]
\addplot[only marks, mark=o] coordinates {(1,1) (2,2) (8,8) (9,9) (1,8) (2,7) (8,2) (9,1)};
\end{axis}
\end{tikzpicture}
\caption{Diferentes tipos de enlace producirían distintas particiones.}
\end{center}

\subsection{Corte del dendrograma}
El dendrograma se corta a una altura determinada para obtener los clusters. La altura puede elegirse mediante:
\begin{itemize}
    \item Criterio visual: buscar un salto grande en las distancias de fusión.
    \item Método de la silueta: probar diferentes cortes y evaluar con silueta.
\end{itemize}

\subsection{Complejidad computacional}
El clustering jerárquico aglomerativo tiene complejidad \(O(n^3)\) en su forma ingenua, o \(O(n^2 \log n)\) con estructuras de datos adecuadas. No es adecuado para grandes conjuntos de datos.

\section{Métricas adicionales para clustering}

Cuando se dispone de etiquetas verdaderas (ground truth), se pueden usar métricas externas.

\subsection{Índice de Rand (Rand Index)}
Mide la similitud entre dos particiones (la real y la obtenida). Para un conjunto de \(n\) puntos, se consideran todos los pares:
\[
RI = \frac{a + b}{\binom{n}{2}}
\]
donde:
\begin{itemize}
    \item \(a\): número de pares que están en el mismo cluster en ambas particiones.
    \item \(b\): número de pares que están en diferentes clusters en ambas particiones.
\end{itemize}
Varía entre 0 y 1, siendo 1 la coincidencia perfecta.

\subsection{Índice de Rand Ajustado (Adjusted Rand Index - ARI)}
Corrige el RI por el azar, esperando que el valor sea 0 para particiones aleatorias:
\[
ARI = \frac{RI - \mathbb{E}[RI]}{\max(RI) - \mathbb{E}[RI]}
\]
Puede tomar valores negativos si la concordancia es peor que la aleatoria. Es la métrica externa más utilizada.

\subsection{Información Mutua (Mutual Information)}
Mide cuánta información comparten dos particiones:
\[
MI(U,V) = \sum_{i=1}^k \sum_{j=1}^l P(i,j) \log \frac{P(i,j)}{P(i)P(j)}
\]
donde \(P(i)\) es la probabilidad de que un punto pertenezca al cluster \(i\) en \(U\), etc.

\subsection{Información Mutua Normalizada (NMI)}
Escala la MI a \([0,1]\) dividiendo por la media de las entropías:
\[
NMI = \frac{MI(U,V)}{\sqrt{H(U)H(V)}}
\]

\subsection{Comparación de métricas}
\begin{table}[h]
\centering
\begin{tabular}{|l|c|c|c|}
\hline
\textbf{Métrica} & \textbf{Rango} & \textbf{Corrige azar} & \textbf{Requiere etiquetas} \\
\hline
Silueta & \([-1,1]\) & No & No \\
Davies-Bouldin & \([0,\infty)\) (menor mejor) & No & No \\
Rand Index & \([0,1]\) & No & Sí \\
ARI & \([-1,1]\) & Sí & Sí \\
NMI & \([0,1]\) & Aproximadamente & Sí \\
\hline
\end{tabular}
\caption{Comparación de métricas de clustering.}
\end{table}

\section{Comparación y cuándo usar cada técnica}

\begin{table}[h]
\centering
\begin{tabular}{|l|l|}
\hline
\textbf{Técnica} & \textbf{Cuándo usarla} \\
\hline
Regresión polinómica & Relación no lineal suave, pocas variables, se necesita interpretabilidad. \\
Regresión robusta (Huber) & Outliers moderados, se desea mantener eficiencia. \\
RANSAC & Gran cantidad de outliers (hasta >50\%), se necesita un modelo de los inliers. \\
SVR & Relaciones no lineales, robustez a outliers, control sobre la banda de error. \\
DBSCAN & Clusters de forma arbitraria, presencia de ruido, no se conoce \(k\). \\
Clustering jerárquico & Se necesita una jerarquía interpretable, pocos datos. \\
\hline
\end{tabular}
\caption{Guía de selección de técnicas complementarias.}
\end{table}

\section{Caso integrado: Análisis de datos de sensores industriales}

\subsection{Problema}
Una planta industrial tiene sensores que registran temperatura, presión y vibración. Se sospecha que ciertas lecturas anómalas corresponden a fallos incipientes. Se dispone de 5000 lecturas sin etiquetar.

\subsection{Objetivos}
\begin{enumerate}
    \item Identificar patrones de comportamiento normal y anomalías.
    \item Agrupar las lecturas en perfiles de operación.
    \item Modelar la relación entre presión y temperatura para predecir valores esperados.
\end{enumerate}

\subsection{Pasos}

\subsubsection{Detección de outliers con RANSAC}
Se ajusta una regresión lineal robusta con RANSAC para modelar temperatura en función de presión. Los puntos que quedan como outliers (no inliers) son candidatos a anomalías.

\subsubsection{Clustering con DBSCAN}
Sobre las variables originales (temperatura, presión, vibración), se aplica DBSCAN. Se elige eps mediante el gráfico de distancia al 5-vecino. DBSCAN identifica varios clusters densos y un conjunto de ruido que coincide en gran medida con los outliers de RANSAC.

\subsubsection{Regresión polinómica por cluster}
Para cada cluster denso, se ajusta una regresión polinómica de grado 2 para modelar la relación temperatura-presión dentro de ese modo de operación. Esto permite establecer límites de control: si una nueva lectura se desvía significativamente de la curva de su cluster, podría indicar un fallo.

\subsubsection{Evaluación}
\begin{itemize}
    \item La silueta promedio de DBSCAN es 0.61, indicando buena separación.
    \item Los outliers detectados por RANSAC y DBSCAN tienen una concordancia del 85\% según el índice de Rand ajustado.
\end{itemize}

\subsection{Conclusión}
La combinación de técnicas complementarias permite una caracterización completa de los datos, identificando modos de operación y anomalías de forma robusta.

\section{Implementación en Python}

\subsection{Regresión polinómica}
\begin{lstlisting}[language=Python]
from sklearn.preprocessing import PolynomialFeatures
from sklearn.linear_model import LinearRegression
from sklearn.pipeline import Pipeline

model = Pipeline([
    ('poly', PolynomialFeatures(degree=3)),
    ('linear', LinearRegression())
])
model.fit(X_train, y_train)
\end{lstlisting}

\subsection{Regresión robusta con Huber}
\begin{lstlisting}[language=Python]
from sklearn.linear_model import HuberRegressor

huber = HuberRegressor(epsilon=1.35)
huber.fit(X_train, y_train)
\end{lstlisting}

\subsection{RANSAC}
\begin{lstlisting}[language=Python]
from sklearn.linear_model import RANSACRegressor

ransac = RANSACRegressor(min_samples=50, residual_threshold=5.0, max_trials=100)
ransac.fit(X_train, y_train)
inlier_mask = ransac.inlier_mask_
\end{lstlisting}

\subsection{SVR}
\begin{lstlisting}[language=Python]
from sklearn.svm import SVR

svr = SVR(kernel='rbf', C=1.0, epsilon=0.1)
svr.fit(X_train, y_train)
\end{lstlisting}

\subsection{DBSCAN con gráfico de distancias}
\begin{lstlisting}[language=Python]
from sklearn.cluster import DBSCAN
from sklearn.neighbors import NearestNeighbors
import numpy as np

# Calcular distancias al k-ésimo vecino (k=minPts)
neigh = NearestNeighbors(n_neighbors=5)
neigh.fit(X)
distances, _ = neigh.kneighbors(X)
k_dist = np.sort(distances[:, -1])

# Graficar para encontrar codo
import matplotlib.pyplot as plt
plt.plot(k_dist)
plt.show()

# Aplicar DBSCAN con eps elegido
dbscan = DBSCAN(eps=0.5, min_samples=5)
labels = dbscan.fit_predict(X)
\end{lstlisting}

\subsection{Clustering jerárquico y dendrograma}
\begin{lstlisting}[language=Python]
from scipy.cluster.hierarchy import dendrogram, linkage, fcluster
import matplotlib.pyplot as plt

Z = linkage(X, method='ward')
plt.figure(figsize=(10,5))
dendrogram(Z)
plt.show()

# Cortar el dendrograma para obtener clusters
clusters = fcluster(Z, t=3, criterion='maxclust')
\end{lstlisting}

\subsection{Métricas de clustering}
\begin{lstlisting}[language=Python]
from sklearn.metrics import adjusted_rand_score, normalized_mutual_info_score

# Si se tienen etiquetas verdaderas y_true
ari = adjusted_rand_score(y_true, labels)
nmi = normalized_mutual_info_score(y_true, labels)
\end{lstlisting}

\section{Conexión con el resto del curso}
\begin{itemize}
    \item La regresión robusta y SVR complementan los modelos de regresión vistos en sesiones anteriores.
    \item DBSCAN y clustering jerárquico amplían las opciones de clustering más allá de K-Means.
    \item Las métricas externas (ARI, NMI) permiten evaluar clustering cuando se dispone de ground truth, algo que no se había tratado.
\end{itemize}

\section{Resumen}
\begin{itemize}
    \item La regresión polinómica modela no linealidades añadiendo términos de mayor grado, pero debe controlarse el overfitting.
    \item La regresión robusta (Huber, RANSAC) es esencial cuando los datos contienen outliers significativos.
    \item SVR extiende SVM a regresión, permitiendo controlar la tolerancia al error mediante \(\varepsilon\) y usar kernels.
    \item DBSCAN es una potente herramienta de clustering basada en densidad, con variantes como OPTICS para densidades variables.
    \item El clustering jerárquico proporciona una visión multinivel de los datos mediante dendrogramas, con diferentes tipos de enlace.
    \item Las métricas externas (ARI, NMI) permiten comparar clusters con etiquetas verdaderas.
    \item La elección de la técnica adecuada depende de la naturaleza de los datos y del objetivo del análisis.
\end{itemize}

\end{document}